\subsection{HUG}
HUG supplies MANO hand hypotheses from observations \cite{hug}. Our conversion
places the native wrist at the Shadow palm position and retargets object-relative
landmarks through Dex Retargeting \cite{dex}. Native and target hands have
different surface geometry, so numerical reconstruction and correct forward
kinematics are necessary but insufficient evidence of grasp preservation.

All 32 saved native exports reconstruct exactly in the recorded 3D audit;
this does not assert all-32 camera-projection equivalence. Detailed surface
measurements on six existing apple, milk-box and mug development cases, together
with three finer repeats, reject the declared strict preservation claim.
Definite lost/added proximity regions and reliable sampled orientation reversals
remain counterexamples. The nine resolutions are repeated measurements of six
grasps, not nine independent physical failures.

The first experimental local fit adds five individual fingertip nonworsening
constraints, a separate strict 10 mm mean-error guard and native-relative facing
targets, including away-facing targets. Final engineering assessment separately
declares a 10 nm arithmetic allowance; original tip errors are approximation
errors, rather than scientific acceptance margins. None of the six fits provides
an eligible improving candidate; each returns its exact original joint/root pose.
Independent source checks pass, while the unchanged surface audit returns zero
PASS and nine FAIL. Pointwise orientation separately has seven FAIL and two
UNRESOLVED results. This is not a repaired mapping or evidence of global
infeasibility. Physical contact and failure-label preservation remain unconfirmed.


A subsequent time-zero diagnosis reproduces all 649 recorded fingertip vectors
exactly and finds full-rank local fingertip-cap Jacobians in all six cases.
Nonzero inward steps satisfying the original caps exist. The unsuccessful
surface-optimization trajectories, rather than contradictory fingertip caps,
explain the first method's refusals; nearby direction-eligibility discontinuities
are observed but are not established as the cause of every refusal.

A separately frozen restoration stage projects each previously rejected raw
pose using cheap forward kinematics before evaluating the unchanged surface
objective. All six restored poses improve that objective by 1.52--38.09\% while
satisfying all 30 individual original-relative fingertip caps in the declared
computed arithmetic and the strict 10 mm mean guard. The projection targets
are tightened inward by 1 nm; no acceptance allowance is added. Fresh terminal
checks bypass the cache, and independent source-URDF versus captured-scene
checks pass. This stage uses 18 full surface evaluations in 49.57 s on the
existing development cases, with no new inference or physical episodes. These
are feasible improving points. Two cheap projections converge and four reach
their 40-iteration budget; final feasibility is checked independently of the
optimizer's status. This is not a full-convergence or surface-preservation claim.

The unchanged independent mapping audit gives zero PASS and nine FAIL; all
nine pointwise orientation results now FAIL. Definite lost/added proximity
regions remain despite the improved fitting cost. For example, the coarse
2 mm milk-box B summed added-area lower bound changes from 142.31 to
245.00 mm$^2$, while its fitting cost falls by 38.09\%. At the coarse 5 mm
band, mug A supplies seven reliable sampled sign-reversal witnesses, compared
with none observed previously. These are conditional anatomical-part area
sums and pose-specific eligible sample counts, not unique contact areas or
comparable failure rates. They show that the average fitting objective and
tip constraints do not certify the complete native grasp. The production
mapper therefore remains unchanged; continuous physical contact and
failure-label preservation are still unconfirmed.
